\chapter{Project Results} 
\label{chap:projectresults}
\setlength{\parindent}{20pt}
%\begin{spacing}{2}
%The goal of the object is to have the speaker, after having the input %signal go through our filters, sound as good as possible to a select set of %judges. It will also be judged on some objective factors, like the gain, %the amount of bass and how flat the frequency response is.
%Our objective is to fulfill these aims to the best of our abilities. %For this we will have to use our knowledge from Linear Circuits A and B.
\paragraph{}
The primary objective of this Integrated Project is to design and synthesize an amplifier for a loudspeaker system consisting of three distinct speaker cones: a woofer, a mid-toner, and a tweeter.
If the primary goals are achieved early, there is the potential to extend the Integrated Project by incorporating the Booming Bass Extension alongside the main amplifier. The amplifier system will be divided into several components: three separate filters tailored to each of the distinct speaker cones, a power supply, and a power amplifier.
\paragraph{}
There are a few specifications and requirements for the loudspeaker; these specifications should be matched to the best of our ability. The target of the amplifier system is to achieve a pass-band gain of 25 times the input voltage. This pass-band has cut-off frequencies of 20 Hz and 40 kHz\cite{janssen-2024,p. 20}. The input signal is expected to have a maximum amplitude of 0.4 volts.  These are just some of the requirements for the amplifier system; a full list of requirements is listed in the chapter \ref{chap:projectactivities} (\nameref{chap:projectactivities}).

The power supply must be designed to provide a positive and a negative voltage. The unloaded voltage should be 20-22 V DC, while maintaining an output voltage of 17-20 V DC (with a maximum of 5\% ripple) when delivering 1.0 A.\cite{janssen-2024, p. 18}

The audio amplifier filters must have a frequency range of 20 Hz to 20 kHz and the different frequency ranges supplement each other well and the transition in frequencies between the different types of filters is as smooth as possible. It is also stated that the amplitude and phase of the different frequency ranges are correct in a way that the total transfer is as flat as possible. \cite{janssen-2024, p.19}

\paragraph{}
If all primary goals are achieved early, an additional challenge is to incorporate a Booming Bass Extension alongside the main amplifier system to further enhance the low-frequency response around 20 Hz. The designing, synthesizing, and implementation of these distinct components will use the knowledge obtained earlier in the Linear Circuits A and B courses.

%\end{spacing}


